% !TEX root = ../main.tex
\section{Diseño del arreglo que almacenará los estudiantes}

El arreglo de registros completos no se carga en memoria principal: vive
en disco, en los tres archivos binarios descritos en la sección anterior.
Lo que sí reside en RAM en todo momento es un arreglo mucho más liviano,
al que llamamos \textbf{mapa de posiciones}: guarda únicamente la llave de
búsqueda, la posición física del registro en disco y su estado, sin el
resto de los datos del estudiante (el funcionamiento completo de este
mapa se detalla en la sección de archivo indexado).

\subsection{Por qué cargar todo en memoria no escala}

\begin{longtable}{p{4.5cm} p{3cm} p{3cm} p{2cm}}
\toprule
Cantidad de estudiantes & Registro completo (3 archivos) & Solo mapas de posiciones (3 mapas) & Ahorro \\
\midrule
\endhead
25\,000 & $\approx$ 6.2 MB & $\approx$ 1.2 MB & 5 veces \\
100\,000 & $\approx$ 25 MB & $\approx$ 4.8 MB & 5 veces \\
1\,000\,000 & $\approx$ 259 MB & $\approx$ 48 MB & 5 veces \\
\bottomrule
\end{longtable}

El cálculo usa 259 bytes por estudiante (suma de los tres structs de la
sección de diseño del registro, sin contar relleno de alineación) contra
16 bytes por entrada de mapa (carné, posición y estado), multiplicado por
tres mapas (por carné, por identificación, por apellido). Aunque un
millón de registros completos (aproximadamente 259 MB) todavía cabría en
la memoria de una máquina moderna, el diseño no depende de esa suposición:
el costo en RAM de los mapas crece mucho más lento que el de los datos
completos, por lo que la arquitectura sigue siendo viable sin importar
cuánto crezca la matrícula de la universidad.
